% Proposed after independent reviews; intentionally not input during cp05 freeze.
\statement{C.37}{Proved}{The canonical word passes every icosahedral generation test}
Let \(\alpha=(-1,f)\) be an involution of the full
\(\widehat{\mathrm{GT}}_0\), and let \(u\in F'\) be its
unique full R1 word, so that \(u=f\alpha(u)\). Then, for
every continuous epimorphism \(q:F\to A_5\) and every
\(\mu\in\widehat{\mathbb Z}^{\times}\),
\[
 \bigl\langle q(x^\mu),q(u^{-1}y^\mu u)\bigr\rangle=A_5.
 \tag{C.37.1}
\]
This is an assertion about each individual finite quotient. It does
not assert generation in the joint quotient consisting of all
nineteen \(A_5\) factors, a relative pro-\(p\) completion,
or the full group \(F\); nor does it identify the third cusp.
No pentagon equation is used.

\begin{proof}
\emph{Two commuting full involutions.}\par
Put \(y'=u^{-1}yu\), and let \(\iota\) interchange
\(x,y\). By C.35, \(\iota(u)=u^{-1}\). Therefore
\[
             v=\operatorname{Inn}(u^{-1})\iota
\]
is an involution interchanging \(x\) and \(y'\).
The full identities
\(\iota\alpha=\operatorname{Inn}(f)\alpha\iota\)
and \(\alpha(u)=f^{-1}u\) give
\[
 \alpha v
 =\operatorname{Inn}\bigl(\alpha(u)^{-1}f^{-1}\bigr)
       \iota\alpha
 =\operatorname{Inn}(u^{-1})\iota\alpha=v\alpha.
 \tag{C.37.2}
\]
Thus \(\tau=\alpha v\) is a full involution. Its action on
integral first homology is minus the interchange matrix, and its
action on mod-two first homology is the regular
\(\mathbb F_2[C_2]\)-module. Both \(v\) and \(\iota\)
preserve the conjugacy class of the third cusp
\(z=(xy)^{-1}\) with exponent one; hence \(\tau\)
preserves this class with exponent minus one.

\emph{A one-section consequence of the full restriction theorem.}\par
Suppose more generally that an involution \(\sigma\) of
\(F\) has the regular mod-two first-homology module.
Write \(E=F\rtimes_\sigma C_2\) and \(k=\mathbb F_2\).
The Hochschild--Serre spectral sequence has just its two fibre
rows, because \(\operatorname{cd}_2(F)=1\). Positive
cohomology of the regular \(k[C_2]\)-module vanishes.
The section makes inflation from \(C_2\) injective, so the
only possible transgression is zero. Consequently
\[
                      \dim_k H^n(E,k)=1\qquad(n\geq2).
 \tag{C.37.3}
\]
Every finite subgroup of \(E\) injects into \(C_2\),
since \(F\) is torsion-free. In particular \(E\) has no
elementary abelian subgroup of rank two. The full profinite
restriction theorem of Kochloukova and Zalesskii, Theorem~2
\cite{CEKochloukovaZalesskii2007}, applies to the open
subgroup \(F\) of finite cohomological dimension. In sufficiently
high degrees it gives a dense restriction embedding into the
product indexed by conjugacy classes of order-two subgroups,
with factors the cohomology of their normalizers. Each factor
contains the one-dimensional cohomology inflated from the
corresponding split \(C_2\). The source in (C.37.3) is finite,
so its dense image is closed and is the whole target. There
can therefore be only one such conjugacy class. There is at
least one, namely the specified section. Conjugacy of sections
by \(E\) is equivalent to conjugacy by \(F\): multiplying
a conjugator by a suitable element of the source section removes
its nontrivial quotient label.

It follows that if \(\sigma\) acts identically on a finite
quotient \(q:F\to Q\) and conjugates a primitive cusp \(z\)
to its inverse, then
\[
                              q(z)=1.
 \tag{C.37.4}
\]
Indeed \(E\to Q\times C_2\) maps the chosen section to
the central element \((1,c)\), and hence maps every section
to this element. Choose a local lift \(r\in E\setminus F\)
that inverts \(z\). Its square belongs to the full procyclic
centralizer \(C_F(z)=\overline{\langle z\rangle}\) from
C.5, and is fixed under conjugation by \(r\). Conjugation by
\(r\) inverts this centralizer, so \(r^2=1\).
Both \(r\) and \(zr\) are involutory sections. Their
images are both \((1,c)\), and their product is \(z\),
proving (C.37.4). No centralizer assertion for the global
section and no torus hypothesis is used here.

\emph{The finite classification used below.}\par
There are \(2280\) ordered generating pairs in \(A_5\),
in nineteen free \(\operatorname{Aut}(A_5)\)-orbits of
size \(120\). A pair's passport here means
\((|a|,|b|,|ab|)\). They divide as follows:
\[
\begin{array}{c|c|l}
\text{generator orders}&\text{number}&\text{additional data}\\ \hline
\{3,5\}&8&\text{both orders occur}\\
\{2,3\}\text{ or }\{2,5\}&6&\text{one generator has order two}\\
(3,3)&1&\text{passport }(3,3,5)\\
(5,5)&4& (5,5,3)\text{ twice},\ (5,5,5),\ (5,5,2)
\end{array}
\tag{C.37.5}
\]
Among the last four, the two generating elements lie in different
\(A_5\)-classes of five-cycles exactly for \((5,5,2)\)
and one of the \((5,5,3)\) orbits.

For each of the six orbits containing an order-two generator,
and for \((3,3,5)\) and the two same-five-class orbits,
a negative peripheral automorphism preserves the quotient
kernel and induces an inner automorphism. This statement uses
the actual conjugacy class of the third cusp, not just generator
orders. Explicitly, for every generating pair \((a,b)\)
in one of these nine orbits, every \(F_0\in A_5\) for
which
\[
 (a^{-1},F_0^{-1}b^{-1}F_0)\text{ generates }A_5,
 \qquad
 a^{-1}F_0^{-1}b^{-1}F_0\sim_{A_5}(ab)^{-1},
 \tag{C.37.6}
\]
gives the same ordered-pair orbit and an inner inducing
automorphism. Also, interchanging the generators preserves
each equal-order orbit; its inducing automorphism is inner
in the \((3,3,5)\) and same-five-class cases, and outer in
the different-five-class cases.

These finite statements are verified by the exhaustive exact
certificate \path{work_CE/code/check_a5_duality_orbits.py}.
It enumerates the sixty projective matrices of
\(\operatorname{SL}_2(\mathbb F_5)\), all one hundred and
twenty automorphisms induced by \(\operatorname{GL}_2(\mathbb F_5)\),
and every possible \(F_0\) in (C.37.6). A separately written
permutation certificate, \path{work_DE/check_review_A5_generation.py},
reproduces the nineteen orbits and all nine tests without importing
the matrix certificate. The class distinction also explains
innerness: inversion preserves each five-cycle class, whereas
the outer automorphism interchanges the two.

We use the elementary subgroup facts that the maximal proper
subgroups of \(A_5\) are the self-normalizing groups
\(A_4,S_3,D_{10}\); no element of \(A_4\) inverts one
of its three-cycles; and the normalizer of a subgroup of odd
order \(3\) or \(5\) is the corresponding dihedral
group. In such an odd dihedral group two involutions commute
only when they are equal. These facts are also independently
checked on all fifty-nine subgroups in the permutation certificate.

\emph{Pairs with distinct generator orders.}\par
Set \(a=q(x)\), \(b=q(y)\), and \(b'=q(y')\).
If their orders are three and five, every conjugate pair
\((a,b')\) still generates: no proper subgroup contains
both orders. This settles the first eight orbits of (C.37.5).

In each of the six cases with an order-two generator, the
kernel-stability assertion (C.37.6) allows us to write the
induced action of \(\alpha\) as \(\operatorname{Inn}(h)\).
Since \(\alpha^2=1\) and \(A_5\) has trivial center,
\(h^2=1\). This action inverts both \(a\) and \(b'\).
If \(M=\langle a,b'\rangle\) is proper, it is one of
\(A_4,S_3,D_{10}\), and \(h\in N_{A_5}(M)=M\).
The \(A_4\) possibility is excluded because \(h\) would
invert a three-cycle. In the odd dihedral possibilities,
\(h\) centralizes the order-two generator and inverts the
odd-order generator, so it equals that order-two generator.

In the finite map
\[
 F\rtimes_\alpha C_2\longrightarrow
 A_5\rtimes_{\operatorname{Inn}(h)}C_2
              \simeq A_5\times C_2,
 \qquad gc\longmapsto(q(g)h,c),
\]
the corresponding one of the full sections \(x\alpha\)
and \(y'\alpha\) therefore has central image \((1,c)\).
The three full section classes represented by
\(\alpha,x\alpha,y'\alpha\) have distinct mod-two
labels. By C.28 their peripheral incidence graph is the
triangle, so the two local reflections at the third cusp
belong to the latter two global classes. If one class has
central finite image, the product of the two reflection
images has order at most two. But the image of the third
cusp in these six passports has odd order three or five.
This contradiction settles all six cases.

\emph{Pairs of equal order.}\par
For either different-five-class orbit, a proper generated
image is cyclic of order five. Write \(b'=a^r\).
The induced involution \(v\) interchanges \(a,b'\),
so \(r^2=1\pmod5\). Different five-cycle classes
require \(r=2\) or \(3\), a contradiction.

It remains to consider \((3,3,5)\) and the two same-five-class
orbits. In these cases both \(\alpha\) and the interchange
\(\iota\) preserve the kernel and induce inner
automorphisms. For a proper \((3,3)\) image the possibility
\(A_4\) is excluded by the same inversion argument as
above; hence the proper image is cyclic of order three.
For a proper \((5,5)\) image it is cyclic of order five.
Again write \(b'=a^r\). The involution \(v\) gives
\(r=1\) or \(-1\).

Suppose first that \(b'=a\). Write the inducing interchange
as \(\operatorname{Inn}(t)\), where \(t\) is a nontrivial
involution of \(A_5\). The inducing inner involution for
\(v\) centralizes \(a\), whose centralizer is an odd
cyclic group. It is therefore identity, and \(q(u)=t\).
Lift the pair \((a,b)\) arbitrarily to matrices
\((A,B)\in\operatorname{SL}_2(\mathbb F_5)^2\), and
let \(T\) lift \(t\). Conjugation by \(T\) interchanges
\(A,B\) up to central signs. A full commutator word is
unchanged by central changes to its inputs, because its two
continuous abelian exponents are zero. Thus, for
\(U=u(A,B)\), full duality gives
\[
                         U\operatorname{Ad}(T)(U)=1.
\]
But \(U=\pm T\), and every lift of a nontrivial
projective involution has square \(-1\): the only
nonidentity involution of \(\operatorname{SL}_2(\mathbb F_5)\)
is \(-1\), as follows directly by diagonalizing a matrix
whose square is one in characteristic five. The displayed
left side is consequently \(-1\), a contradiction.

Finally suppose \(b'=a^{-1}\). The induced inner
involutions \(\alpha\) and \(v\) both invert \(a\)
and lie in its odd dihedral normalizer. They commute by
(C.37.2); their conjugators commute too because \(A_5\)
is center-free. They therefore coincide. The full involution
\(\tau=\alpha v\) acts identically on this finite
quotient. Its mod-two action is regular, and it inverts the
primitive third cusp up to conjugacy. Equation (C.37.4)
would make the nontrivial third-cusp image trivial. This final
contradiction proves generation for all nineteen orbits.

We have proved (C.37.1) for \(\mu=1\). For arbitrary
\(\mu\in\widehat{\mathbb Z}^{\times}\), each generator's
unit power generates the same finite cyclic subgroup.
The generated subgroup is therefore unchanged, proving
the stated unit-powered version as well.
\end{proof}

The finite certificate also enumerates every proper-image
candidate satisfying the necessary binary duality condition.
There are \(342\): \(108\) have no required inner inverting
involution, \(208\) have the central-section contradiction,
and \(26\) have the one-section contradiction. This count
checks coverage of the cases above. It makes no assertion that
an independently chosen finite evaluation table is realized by
a full canonical word.
